% Standalone polished architecture template. Compile this file with pdfLaTeX.
% No external input files, images, bibliography database, or shell escape needed.
% Adapted from Rich Hilliard's templates, copyright 2012--2014, CC BY 3.0.
% This 2026 working edition changes layout, wording, fields, and edition guidance.
% Preserve attribution and identify further changes when reusing this adaptation.
\documentclass[11pt,letterpaper]{article}
% Polished architecture-template style, working edition 1.0, 2026-09-06.
% Template content is adapted from Rich Hilliard's CC BY 3.0 templates.
\RequirePackage[T1]{fontenc}
\RequirePackage[utf8]{inputenc}
\RequirePackage{mathpazo}
\RequirePackage[scaled=0.92]{helvet}
\RequirePackage{microtype}
\RequirePackage[letterpaper,left=0.77in,right=0.77in,top=0.73in,bottom=0.73in,headheight=21pt,headsep=18pt,footskip=25pt]{geometry}
\RequirePackage[table]{xcolor}
\RequirePackage{array,tabularx,booktabs}
\RequirePackage{enumitem}
\RequirePackage[skins,breakable]{tcolorbox}
\RequirePackage{titlesec,fancyhdr,lastpage,needspace}
\RequirePackage{xurl}
\RequirePackage{hyperref}
\definecolor{Navy}{HTML}{193443}
\definecolor{Teal}{HTML}{137980}
\definecolor{Ink}{HTML}{293D47}
\definecolor{Muted}{HTML}{62737D}
\definecolor{Line}{HTML}{D6E2E6}
\definecolor{Pale}{HTML}{F1F6F7}
\definecolor{Tint}{HTML}{EBF5F3}
\color{Ink}
\hypersetup{colorlinks=true,linkcolor=Teal,urlcolor=Teal,citecolor=Teal}
\urlstyle{same}
\setlength{\parindent}{0pt}
\setlength{\parskip}{6pt}
\setlength{\emergencystretch}{1.5em}
\linespread{1.035}
\raggedbottom
\renewcommand{\arraystretch}{1.2}
\setlength{\tabcolsep}{5pt}
\newcolumntype{L}[1]{>{\raggedright\arraybackslash}p{#1}}
\newcolumntype{Y}{>{\raggedright\arraybackslash}X}
\setlist[itemize]{leftmargin=1.25em,itemsep=3pt,topsep=4pt,parsep=0pt}
\setlist[enumerate]{leftmargin=1.55em,itemsep=4pt,topsep=4pt,parsep=0pt}
\titleformat{\section}{\sffamily\bfseries\Large\color{Navy}}{\textcolor{Teal}{\thesection}}{0.65em}{}
\titleformat{\subsection}{\sffamily\bfseries\normalsize\color{Navy}}{}{0pt}{}
\titlespacing*{\section}{0pt}{0pt}{9pt}
\titlespacing*{\subsection}{0pt}{11pt}{4pt}
\providecommand{\TemplateCode}{ARCHITECTURE TEMPLATE}
\providecommand{\TemplateShortTitle}{WORKING EDITION}
\newcommand{\TemplatePageBreak}{\clearpage}
\newcommand{\FillIn}[1]{\textcolor{Teal}{\sffamily\itshape[#1]}}
\newcommand{\RecordID}[1]{\textsf{\textbf{#1}}}
\newcommand{\TableHead}[1]{\textcolor{Navy}{\sffamily\bfseries#1}}
\newcommand{\Lead}[1]{\textcolor{Muted}{#1}\par\vspace{3pt}}
\newcommand{\Repeat}[1]{\textcolor{Teal}{\sffamily\small\textbf{REPEAT}\enspace #1}\par}
\newcommand{\Guide}[1]{\par{\small\textcolor{Muted}{\textbf{Authoring guidance.} #1}}\par}
\newtcolorbox{TemplateNote}[1]{enhanced,colback=Tint,colframe=Tint,
  coltitle=Navy,title={#1},fonttitle=\sffamily\bfseries,boxrule=0pt,arc=1pt,
  borderline west={2pt}{0pt}{Teal},left=10pt,right=10pt,top=6pt,bottom=7pt,
  before skip=9pt,after skip=9pt}
\newtcolorbox{TemplateRecord}[1]{colback=white,colframe=Line,colbacktitle=Pale,
  coltitle=Navy,title={#1},fonttitle=\sffamily\bfseries\small,boxrule=0.5pt,
  arc=1pt,left=9pt,right=9pt,top=6pt,bottom=6pt,before skip=8pt,after skip=8pt}
\pagestyle{fancy}
\fancyhf{}
\fancyhead[L]{\sffamily\scriptsize\color{Muted}\TemplateShortTitle}
\fancyhead[R]{\sffamily\scriptsize\color{Muted}EDITABLE WORKING TEMPLATE}
\fancyfoot[L]{\sffamily\scriptsize\color{Muted}\TemplateCode\enspace\textbullet\enspace EDITION 1.0}
\fancyfoot[R]{\sffamily\scriptsize\color{Muted}\thepage\enspace /\enspace\pageref*{LastPage}}
\renewcommand{\headrulewidth}{0.4pt}
\renewcommand{\footrulewidth}{0pt}
\renewcommand{\headrule}{\hbox to\headwidth{\color{Line}\leaders\hrule height \headrulewidth\hfill}}
\newcommand{\TemplateTitle}[3]{%
  {\sffamily\small\bfseries\color{Teal}ARCHITECTURE DOCUMENTATION\enspace /\enspace #1\par}
  \vspace{14pt}
  {\sffamily\bfseries\fontsize{29}{33}\selectfont\color{Navy}#2\par}
  \vspace{10pt}
  {\sffamily\large\color{Muted}#3\par}
  \vspace{11pt}
  {\sffamily\small\color{Muted}Polished working edition 1.0\enspace\textbullet\enspace 6 September 2026\par}
  \vspace{13pt}{\color{Teal}\hrule height 1.2pt}\vspace{9pt}}
\newcommand{\TemplateAttribution}{%
  \textbf{Attribution.} Adapted from the \emph{Architecture Description Template}
  and \emph{Architecture Viewpoint Template} by Rich Hilliard,
  \copyright\ 2012--2014, supplied under
  \href{https://creativecommons.org/licenses/by/3.0/}{Creative Commons Attribution 3.0 Unported}.
  This adaptation revises layout, wording, fields, dependencies, and edition guidance;
  it does not imply endorsement by the original author or the standards bodies.}
\newcommand{\EditionNote}{%
  The supplied templates were written for ISO/IEC/IEEE 42010:2011.
  This adaptation uses 2022 terminology and adds records for perspectives,
  aspects, and view components. It is an authoring aid, not an official template
  or a complete normative checklist. Assess the completed work against the
  published edition before making a conformance claim.}
\newcommand{\TemplateSources}{%
  {\small
  \begin{enumerate}[label={[\arabic*]},leftmargin=1.6em,itemsep=6pt]
  \item Rich Hilliard. \emph{Architecture Description Template} and
  \emph{Architecture Viewpoint Template}, revision 2.2 (2014), with the supplied
  viewpoint section marked version 2.1b.
  \href{http://www.iso-architecture.org/42010/templates/}{Original source location}.
  The attached source files are the adaptation baseline.
  \item ISO. \emph{ISO/IEC/IEEE 42010:2022: Software, systems and enterprise---Architecture description}.
  \href{https://www.iso.org/standard/74393.html}{Official catalogue and scope}.
  \item ISO/IEC/IEEE. \emph{42010:2022, official preview}.
  \href{https://webstore.ansi.org/preview-pages/iso/preview_iso\%2Biec\%2Bieee\%2B42010-2022.pdf}{Contents and edition changes, hosted by ANSI}.
  \item Creative Commons. \href{https://creativecommons.org/licenses/by/3.0/}{Attribution 3.0 Unported license}.
  \end{enumerate}}}
\renewcommand{\TemplateCode}{VP-GUIDE}
\renewcommand{\TemplateShortTitle}{ARCHITECTURE VIEWPOINT GUIDE}
\title{Architecture Viewpoint\\Template}
\author{Adapted from templates by Rich Hilliard}
\date{6 September 2026}
\hypersetup{pdftitle={Architecture Viewpoint Template},pdfauthor={Adapted from templates by Rich Hilliard},pdfsubject={Editable working template with 2022 edition guidance}}
\makeatletter
\renewcommand{\maketitle}{\TemplateTitle{VIEWPOINT GUIDE}{\@title}{Authoring guide and specification outline}}
\makeatother
\begin{document}
\maketitle
\Lead{Guidance for selecting, defining, and maintaining a reusable architecture
viewpoint, with a compact outline and an illustrative application.}

\begin{TemplateNote}{Edition and use}
\EditionNote
\end{TemplateNote}

\subsection{Use the document that matches the work}
{\small
\begin{tabularx}{\linewidth}{@{}L{1.7in}Y@{}}
\toprule
\TableHead{Document or source} & \TableHead{Purpose}\\
\midrule
Architecture description template & Organize one AD, its participants, concerns,
viewpoints, views, decisions, and evidence.\\
This viewpoint guide & Understand the specification outline, choose conventions,
and plan how the viewpoint will be used.\\
Viewpoint specification worksheet & Complete repeatable fields for a particular
viewpoint, its model kinds, methods, and sources.\\
Reusable viewpoint section & Insert the worksheet body into an existing LaTeX
document using the supplied style support.\\
\bottomrule
\end{tabularx}}

\subsection{What a useful viewpoint provides}
A viewpoint identifies the questions its views are intended to address and
the conventions used to construct and interpret them. It should help an author
select representations and help a reader understand their meaning and limits.
The actual architecture content belongs in the resulting views.

\textbf{Authoring sequence:} establish audience and concerns; define
representations; specify methods; identify correspondences; provide examples
and sources; review the specification against a representative use case.

\vfill
{\footnotesize\color{Muted}\TemplateAttribution\par}

\TemplatePageBreak
\section{Viewpoint specification outline}
\Lead{Use the outline to plan the specification, then complete the companion worksheet.}

{\small
\begin{tabularx}{\linewidth}{@{}L{1.48in}Y@{}}
\toprule
\TableHead{Information group} & \TableHead{Content to develop}\\
\midrule
Identity and overview & Name, synonyms, identifier, version, purpose, author,
steward, applicability, scope, and limitations.\\
Concerns and audience & Concern statements, typical stakeholder roles,
perspective groupings, aspects, and coverage rationale.\\
Model kinds & A catalogue plus a specification of each kind's vocabulary,
notation, semantics, constraints, and adopted sources.\\
Other component formats & Content structure, source information, interpretation
legends, and limitations for non-model-based material.\\
View methods & Construction, interpretation, analysis, and other relevant
operations; inputs, procedure, outputs, and evidence.\\
Correspondence methods & Relationship semantics, permitted participants,
establishment and checking procedures, and violation handling.\\
Examples and notes & Demonstrations, scope boundaries, interpretation advice,
complementary viewpoints, and known limitations.\\
Sources and stewardship & Prior art, versioned references, authorship, change
history, review result, and release status.\\
\bottomrule
\end{tabularx}}

\subsection{Separate reusable conventions from their application}
{\small
\begin{tabularx}{\linewidth}{@{}L{1.0in}Y Y@{}}
\toprule
\TableHead{Topic} & \TableHead{Viewpoint specification} & \TableHead{Project AD and views}\\
\midrule
Stakeholders & Typical roles and their expected questions. & Actual participants or
applicable classes and their recorded concerns.\\
Model kind & Vocabulary, permitted relationships, notation, and interpretation. &
The particular architecture elements and relationships represented.\\
Method & Inputs, procedure, result meaning, and limitations. & Execution baseline,
evidence, findings, and unresolved results.\\
Correspondence & Meaning of a relationship and how to establish or check it. &
Actual participating items and the recorded result of any assessment.\\
\bottomrule
\end{tabularx}}

\Guide{A viewpoint can be reused across projects and can govern more than one
view. Do not treat the specification as a single drawing or require one drawing
to carry every concern. Keep its intended coverage explicit.}

\TemplatePageBreak
\section{Define representations and methods}
\Lead{Specify meaning first. Styling then makes that meaning easier to read.}

\subsection{Three complementary ways to describe a model kind}
\begin{enumerate}
\item \textbf{Adopt a language.} Name the language or notation, its edition,
the constructs used, and any restrictions or extensions. Identify the source
that determines how those constructs are interpreted.
\item \textbf{Define a metamodel.} Describe element types, properties,
relationship types, and constraints. Explain the distinction between a rule
for the representation and a constraint on the architecture being represented.
\item \textbf{Provide a structured template.} Define the fields, their types,
permitted values, references, and interpretation. Give an example of a correctly
completed instance.
\end{enumerate}

\subsection{Minimum useful convention record for authors}
{\small
\begin{tabularx}{\linewidth}{@{}L{1.25in}Y@{}}
\toprule
\TableHead{Question} & \TableHead{Make the answer explicit}\\
\midrule
What may appear? & Element and relationship vocabulary; IDs; attributes; units;
source and scope conventions.\\
What does it mean? & Semantics of direction, containment, connections, labels,
status, and any visual encoding.\\
What is permitted? & Well-formedness rules, constraints, assumptions, and
versioned language or profile references.\\
How is it read? & Legend, reading order where relevant, abstraction level,
interpretation boundaries, and a worked example.\\
How is it assessed? & Inputs, procedure, criteria, outputs, evidence, and
handling of incomplete or contradictory information.\\
\bottomrule
\end{tabularx}}

\subsection{View-method specification}
For each method, record an identifier, purpose, applicable concerns and model
kinds, input baselines, preconditions, steps, outputs, and limitations. Distinguish
a method definition from the evidence produced when it is applied.

\textbf{Construction} explains how to produce the view.
\textbf{Interpretation} explains how to understand it.
\textbf{Analysis} explains how to reason from it or assess it.
Other operations, including transformation or implementation guidance, can be
included where the viewpoint supports those uses.

\begin{TemplateNote}{Support non-model-based content deliberately}
\small
A narrative explanation or evidence table can be a useful view component.
Define its format, sources, and legend. A diagram type by itself does not specify
the complete vocabulary, meaning, and assessment approach of a model kind.
\end{TemplateNote}

\TemplatePageBreak
\section{Illustrative viewpoint: service continuity}
\Lead{This example demonstrates how a specification can connect a concern to
representations and assessment methods. It is not a validated architecture.}

{\small
\begin{tabularx}{\linewidth}{@{}L{1.38in}Y@{}}
\toprule
\TableHead{Specification field} & \TableHead{Illustrative content}\\
\midrule
Viewpoint name & Service continuity; identifier VP-CONTINUITY.\\
Audience & Service operators, developers, and service owners.\\
Concern & What happens to a service when an external dependency is unavailable?\\
Perspective and aspects & Operational continuity; structure, behavior, and timing.\\
Applicability & Services with identifiable dependencies and documented failure scenarios.\\
Model kinds & Dependency graph and failure-scenario table.\\
Other components & Narrative assumptions, evidence summaries, and unresolved-issue records.\\
Limitations & A dependency graph alone cannot establish recovery time or successful failover.\\
\bottomrule
\end{tabularx}}

\subsection{Representation conventions}
\textbf{Dependency graph:} a node represents an identified service or external
dependency. A directed relation from A to B means A relies on B for a stated
operation. Each relation identifies that operation and the relevant baseline.

\textbf{Failure-scenario table:} each record identifies the failed dependency,
affected operation, triggering condition, proposed response, recovery criterion,
and evidence or unresolved status. Any timing criterion has a declared source.

\subsection{Methods and evidence}
\begin{enumerate}
\item Construct the dependency representation from declared interfaces and
identify the evidence supporting each dependency.
\item Record failure scenarios for the concerns selected by the stakeholders.
Identify assumptions about timeouts, retry behavior, and recovery.
\item Compare scenario participants with the dependency representation.
Record missing, inconsistent, or ambiguous identifiers.
\item Evaluate the behavior using an appropriate review, simulation, or test.
Retain the baseline, observations, limitations, and unresolved findings.
\end{enumerate}

\begin{TemplateNote}{Correspondence example}
\small
The dependency named in a failure scenario refers to the same identified
dependency in the graph. A checking method compares identifiers and the
affected operation across both representations. The resulting AD records the
relationship and the outcome of that check separately.
\end{TemplateNote}

\TemplatePageBreak
\section{Review, edition guidance and sources}
\Lead{Release a usable specification with a known baseline and explicit limits.}

\subsection{Review the specification by using it}
Ask a representative author to construct a small view using the specification,
then ask a representative reader to interpret it. Record ambiguous notation,
missing inputs, unanswered concerns, and methods that cannot be applied as written.

{\small
\begin{tabularx}{\linewidth}{@{}L{1.32in}Y@{}}
\toprule
\TableHead{Review record} & \TableHead{Entry}\\
\midrule
Specification baseline & \FillIn{Viewpoint identifier, version, sources, and reviewed artifact}\\
Participants and use case & \FillIn{Author and reader roles; example scope and concern set}\\
Findings and disposition & \FillIn{Ambiguities, gaps, changes, owners, and evidence}\\
Release and revisit & \FillIn{Release status, steward, date, and triggers for review}\\
\bottomrule
\end{tabularx}}

\subsection{How this edition treats the earlier outline}
The supplied files target the 2011 edition. This adaptation retains their
identity, concerns, stakeholders, model kinds, operations, examples, notes,
and sources, while adding 2022-oriented records and removing unsupported
claims that completing the template establishes conformance.

The 2022 publication uses \emph{entity of interest}, explicitly addresses
perspectives and aspects, and accommodates view components that are not
model-based. Its contents locate viewpoint and model-kind specifications in
Clause 8; older references to Clause 7 should not be carried forward as 2022
clause references. See the official edition preview below.

\subsection{References and attribution}
\TemplateSources
{\footnotesize\color{Muted}\TemplateAttribution\par
\textbf{Working-edition changes:} consistent typography and editable outlines;
self-contained sources; explicit separation of guide, worksheet, and reusable
section; new coverage, interpretation, method, and evidence fields.
The published standard remains the basis for assessing any conformance claim.\par}

\end{document}
